\section{Исследование параметров фильтрации и усреднения трассы}

В данном разделе исследуется влияние полосы пропускания фильтра промежуточной частоты (\textit{RBW}), полосы пропускания видеофильтра (\textit{VBW}) и режима усреднения трассы по мощности (\textit{Power Avg}) на разрешающую способность, уровень шумов и время развёртки анализатора спектра при фиксированных частоте сигнала $f_0 = 1{,}0\text{ ГГц}$ и полосе обзора $\text{Span} = 10\text{ МГц}$.

\subsection{Влияние уменьшения полосы пропускания фильтра ПЧ (RBW = 300 Гц)}

Для оценки предельной избирательности полоса фильтра промежуточной частоты была принудительно уменьшена до минимального рекомендуемого значения $\text{RBW} = 300\text{ Гц}$ при сохранении исходного значения видеофильтра $\text{VBW}$. Полученная спектрограмма приведена на рис.~\ref{fig:rbw_300hz}.

\begin{figure}[H]
    \centering
    \includegraphics[width=0.8\linewidth]{fig6.png}
    \caption{Спектрограмма при минимальной полосе фильтра ПЧ ($\text{RBW} = 300\text{ Гц}$)}
    \label{fig:rbw_300hz}
\end{figure}

При уменьшении полосы $\text{RBW}$ форма спектрального отклика становится существенно уже и острее, приближаясь к теоретическому виду линейчатого спектра гармонического колебания. 

Одновременно с этим наблюдается выраженное снижение среднего уровня шумовой дорожки ($\text{DANL}$). Это обусловлено тем, что интегральная мощность теплового шума на входе детектора прямо пропорциональна ширине полосы тракта промежуточной частоты: $P_{\text{ш}} = G_0 \cdot \text{RBW}$.
Главным недостатком сужения полосы $\text{RBW}$ выступает существенное увеличение времени цикла развёртки (\textit{Sweep Time}).
\subsection{Влияние изменения полосы видеофильтра (VBW)}

Для проведения исследования полоса пропускания фильтра ПЧ была возвращена к значению $\text{RBW} = 100\text{ кГц}$. Это позволило значительно ускорить процесс развёртки и отчетливо зафиксировать влияние полосы видеофильтра $\text{VBW}$. В ходе эксперимента сравнивались два предельных режима: сужение полосы до $\text{VBW} = 10\text{ Гц}$ и расширение до $\text{VBW} = 3\text{ МГц}$. Результаты измерений представлены на рис.~\ref{fig:vbw_10hz} и рис.~\ref{fig:vbw_3mhz}.

\begin{figure}[H]
    \centering
    \begin{minipage}[b]{0.48\textwidth}
        \centering
        \includegraphics[width=\textwidth]{fig7.png}
        \caption{Полоса видеофильтра заужена ($\text{VBW} = 10\text{ Гц}$)}
        \label{fig:vbw_10hz}
    \end{minipage}
    \hfill
    \begin{minipage}[b]{0.48\textwidth}
        \centering
        \includegraphics[width=\textwidth]{fig8.png}
        \caption{Полоса видеофильтра расширена ($\text{VBW} = 3\text{ МГц}$)}
        \label{fig:vbw_3mhz}
    \end{minipage}
\end{figure}

Видеофильтр представляет собой низкочастотный фильтр (ФНЧ), включённый после детектора огибающей. Он предназначен для сглаживания шумов и не влияет на частотную избирательность прибора по ВЧ-входу. 

При установке $\text{RBW} = 300\text{ Гц}$ вариация $\text{VBW}$ практически не давала видимого эффекта. Это объясняется тем, что спектр шумов на выходе амплитудного детектора физически ограничен полосой предшествующего тракта ПЧ (то есть не превышает $300\text{ Гц}$).
После расширения полосы до $\text{RBW} = 100\text{ кГц}$ на детектор начал поступать широкий спектр шумов с высокой дисперсией. При расширении видеофильтра до $\text{VBW} = 3\text{ МГц}$ (когда $\text{VBW} > \text{RBW}$) все быстрые случайные флуктуации шума беспрепятственно проходят на экран, образуя широкую «изрезанную» шумовую дорожку с большим разбросом значений. 

При сужении видеофильтра до $\text{VBW} = 10\text{ Гц}$ (режим $\text{VBW} \ll \text{RBW}$) фильтр эффективно усредняет высокочастотные шумовые броски, превращая шумовую дорожку в гладкую непрерывную линию. Однако из-за большой постоянной времени низкочастотного фильтра быстродействие прибора заметно падает, требуя увеличения времени цикла развёртки.

\subsection{Усреднение трассы мощности}

Для демонстрации цифровых методов подавления шумов были выставлены одинаковые полосы фильтрации $\text{RBW} = \text{VBW} = 3\text{ кГц}$. В меню настроек трассы был активирован режим усреднения по мощности (\textit{Power Avg}) со значением счётчика $N = 1000$ отсчётов. Результат усреднения приведён на рис.~\ref{fig:trace_avg_1000}.

\begin{figure}[H]
    \centering
    \includegraphics[width=0.8\linewidth]{fig9.png}
    \caption{Спектрограмма в режиме усреднения по 1000 отсчётам ($\text{Power Avg}$)}
    \label{fig:trace_avg_1000}
\end{figure}

Накопление и усреднение спектра по $1000$ проходам развёртки приводит к сглаживанию случайных флуктуаций шумов без расширения полосы пропускания аналоговых трактов. За счёт усреднения спектр сигнала преобразуется в стабильный пик, что повышает достоверность измерений амплитуды гармоники маркером.

Полное время формирования стабильной картины спектра возрастает прямо пропорционально количеству усредняемых реализаций $N$. Из-за высокой инерционности такой режим пригоден только для анализа стационарных колебаний и может приводить к потере информации о быстрых импульсных помехах или нестабильных сигналах.
